% ===== §4 =====
\section{The First Impossibility: The Non-Structurability of Value}
\label{sec:structurability}

\noindent
The four impossibilities descend from the most superficial precondition of the anatomy to the most fundamental, and it is worth being explicit at the outset about why this first one comes first, since the order is an argument in itself. The later impossibilities all concern a function already written: whether it can be held fixed, whether its values share a scale, whether a stable subject owns it. They presuppose, that is, that the values at stake have already been got onto the page as terms, and they ask what goes wrong thereafter. This first impossibility strikes earlier, at the writing itself. It asks whether certain of the values any honest account of intimate relation must weigh can be set down as terms at all, and it answers that some cannot, so that the objective is not merely hard to fix or to compare but impossible to complete. A value that cannot be written into $J$ is not a value the optimization prices badly; it is a value the optimization never reaches, excluded not from the maximum but from the statement of the problem. If this is right, the optimizing ambition fails at the first step, before the later impossibilities are even in view, and the completeness on which the whole ambition silently rests is shown to have been unavailable from the start.

\subsection*{The completeness demand, made precise}

The claim must be sharpened before it can be defended, because in a loose form it is merely the tired complaint that love is too rich for numbers, and that complaint the economist can absorb without a scratch. His reply, rehearsed already in the introduction, is that a value which seems to resist the function resists only because the function is written too poorly, and that a richer specification, with the goods of intimacy and the long horizon and the pleasures of the bond all entered as arguments, would take up whatever a first pass left out. The reply has real force, and it has a distinguished lineage: the program of extending utility maximization into the household, into marriage, fertility, and the allocation of affection, was carried out in detail and with great ingenuity, and it proceeded exactly by folding ever more of intimate life into the arguments of a utility function \citep{robbins1932nature}. To meet it one must state precisely what it demands. It demands \emph{completeness}: that for every value $v$ bearing on the relation there exist a term $t_v$, a representation of $v$ as an argument of $J$ such that the candidate states of the relation are ordered by how much of $v$ they deliver. Completeness is not an incidental wish of the program but its condition of adequacy, for an objective that omits a value that genuinely bears on the relation is an objective for some thinner thing, and to maximize it is to optimize a caricature. The economist needs completeness, and the question is whether he can have it.

He can have it for a great many values, and the paper concedes this freely, since the concession sharpens the point of attack. Comfort, shared leisure, the division of labor in a household, even many of the tendernesses of a life together admit of proxies stable enough to serve as terms, and where they do the optimizing frame is not obviously ill-posed but merely, as the later parts will argue, ethically and practically treacherous. The first impossibility does not rest on the mass of ordinary goods that can be proxied. It rests on the claim that \emph{at least one} value the objective must include cannot be, and that its exclusion is fatal to completeness however well the rest is done. The paper meets the demand, therefore, not at its weakest point but at its hardest: with the one value whose inclusion the economist most needs and least can secure, the psychoanalytic drive.

\subsection*{The drive as the value that cannot be written}

That the drive must be in any complete objective for intimate relation is not a point the economist can concede away, for the drive is not an ornament at the edge of the bond but among its springs. What draws two people and holds them, what makes a particular other the one whose presence is not substitutable by another of equal merit, what gives desire its insistence and its opacity even to the one who feels it, belongs to the economy of the drive; an objective for a love that left it out would be an objective for companionship or for mutual advantage, some real thing but not the thing in question. So the drive must be a term. And the drive is precisely what admits of no term, for reasons that lie in its structure and not in our instruments.

The drive, in the register the series has used throughout, arises at a coupling of the Real and the Symbolic and reduces to neither \citep{lacan1977xi,lacan1992vii}. It is not a preference. A preference is a settled disposition to rank options, and a ranking is exactly what a term records; if the drive were a preference it would pose no problem, and could be written down like any other. But the drive does not rank the options in the feasible set, because what it circles is not among them. It orbits an object that is not a positive object at all but a vacancy, the place of what is constitutively missing, and it draws its force from the orbiting rather than from any attainable state that might be reached and possessed. This has a consequence that is fatal to representation as a term. A term $t_v$ must assign to each candidate a magnitude, so that candidates can be ordered by it, more here, less there, monotone in the good; and the assignment must be such that getting more of the good is, other things equal, getting nearer the objective's maximum. The drive violates this at its core, for it is extinguished by attainment where attainment occurs: the state that would count, on any proxy, as the fullest delivery of the good is the very state in which the drive, having reached its ostensible object, finds that the object was never it, and moves on. The mapping a term requires, from more of the good to nearer the maximum, is not merely inaccurate for the drive; it runs backwards, and no monotone magnitude can be laid over a structure whose satisfaction is its own undoing.

\begin{casebox}{the ledger of a marriage}
Consider the attempt made in earnest rather than in theory. A couple, urged toward the frame by a culture that prizes it, resolve to reason about their marriage as an optimization: to enumerate what each values, to weigh the yields against the costs, and to act on the balance. Much of the enumeration succeeds. They can list the division of housework, the comfort of a shared home, the pleasures of company, and can even, with effort, assign these rough weights and speak of trade-offs among them without absurdity. Then one of them tries to enter the thing that actually binds them to this person and not another, equal on every listed count, and the enumeration fails. It is not that the value is large and hard to weigh. It is that when they try to say what it is a quantity \emph{of}, every candidate answer, security, admiration, pleasure, familiarity, turns out to name something that a different person could in principle supply as well or better, and so to name not the bond but a proxy the bond exceeds. The remainder that resists every proxy, the sheer this-one-ness that is why the ledger is being kept at all, is precisely what will not go into a column, because it is not a quantity of anything and answers to no ``more'' or ``less.'' The ledger can be completed only by leaving out the reason for keeping it.
\end{casebox}

To force the drive into a term one would therefore have to replace it with a proxy, some measurable satisfaction standing where it stood; and here the decisive move of the economist's program can be named and refused. The program's most powerful device is revealed preference: the resolve to sidestep the inner life entirely and read value off behavior, defining what a person values as whatever their choices maximize, so that the drive, whatever it is in itself, is captured by the pattern of what the person is observed to choose. This is ingenious and, for a wide class of goods, unobjectionable. Against the drive it fails, and its failure is instructive. Revealed preference assumes that choice reveals a stable underlying ordering, that the pattern of choices is the trace of a preference the choices express. But the drive produces exactly the pattern revealed preference cannot parse: the pursuit that abandons its object on attainment, the choice that undoes itself, the wanting whose satisfaction is its extinction. Faced with this, revealed preference must either declare the behavior irrational, thereby excluding from the account the very spring of the bond it was meant to capture, or fit an ordering to the choices post hoc, thereby recording as a stable preference a movement whose whole character is that it is not one. In neither case has it represented the drive. It has, in the first, thrown the drive out as noise, and in the second, replaced it with a fiction of stability that is a different thing wearing its name. The proxy is not a poorer measurement of the drive. It is the substitution of a positive object for a vacancy, which is exactly the operation that empties the drive of what makes it the drive.

\subsection*{The category of the failure, and its consequence}

The nature of the failure must be fixed exactly, because everything downstream depends on it. The failure is not epistemic. It is not that the drive is hard to quantify, that our instruments are coarse and our surveys noisy, and that a finer psychometrics would in time get its number. That reading leaves the optimizing frame wholly intact, conceding only that the measurements are poor; and it is the reading the paper most needs to refuse, because it is the one that keeps the door open to ``do the sums better.'' The claim is stronger and stranger: that there is here no number to be measured well or badly, that the mapping from the drive to a magnitude fails not in its accuracy but in its existence, because the drive is not the kind of entity of which there is a quantity to be had. This is the same distinction the anatomy drew between a problem that is merely hard and a problem that is ill-posed, and the drive falls on the second side of it not by degree but by kind.

\begin{claim}[Non-structurability]
A value is \emph{non-structurable} if its manner of being precludes its representation as a term in an objective function, that is, precludes any assignment to the candidate states of a monotone magnitude by which they could be ordered. The psychoanalytic drive is non-structurable: arising at a coupling of the Real and the Symbolic, circling a constitutive vacancy rather than a positive object, and extinguished by attainment where attainment occurs, it admits no monotone magnitude over the candidates and so no term. Its exclusion from the objective is therefore not a defect of measurement, remediable by finer instruments, but a matter of category. It cannot enter the statement of the problem.
\end{claim}

The consequence is fatal to the economist's reply on its own terms, and it is worth drawing the dilemma out to its end so that no escape is left standing. The reply promised to rescue optimization by completing the function; completeness for intimate relation requires a term for the drive; the drive admits no term. So the function must either omit the drive, in which case it is incomplete and the reply has not been made, the objective being one for some thinner relation that leaves out a spring of this one, or include a proxy for the drive, in which case it has put a positive object where a vacancy was and is again not an objective for this relation but for a substitute, so that the optimization, however impeccably posed thereafter, maximizes the wrong thing at the ground floor. There is no third path on which the drive is both admitted and itself, because admission as a term just is the imposition of a magnitude, and the drive admits none. The first precondition of the anatomy, that every value at stake be representable in the objective, therefore fails, and it fails before the questions of fixity, comparability, and the optimizer can be raised, since those concern a function already written and this concerns the writing. One may read the three impossibilities that follow as showing that even the incomplete function one could write, over such values as do admit of terms, would still furnish no well-posed problem. But this first has already shown that the function was condemned to be incomplete, and that the completeness on which the whole optimizing ambition depended was never there to be had.
